\scanpage{4}
\mspage{( 3 )}

\textit{Alors toute $\mathrm A$-algèbre intègre qui est localisée d’une $\mathrm A$-algèbre finie (i. e. $\mathrm A$-module de type fini) $\mathrm C$ est analytiquement réduite.}
\margin{\handwritten{$C\supset A$ !!!}}
% A curved pencil line leads from C to the upper note.

(Pour appliquer le lemme, on remarque que \tikz[baseline=(x.base)]{\node[draw,ellipse,inner sep=1pt] (x) {$h_1$};} et $h_2$ sont stables par localisation.)
\margin{\handwritten{faux}}
% The line from faux ends at the oval round h_1.

De l’existence de lemmes de normalisation [pour les séries convergentes, \textit{voir} (\textsuperscript{5}) par exemple], on déduit :

\textsc{Corollaire 1} (Chevalley). --- \textit{Toute localité intègre est analytiquement réduite.}

\textsc{Corollaire 2.} --- \textit{Soient $k$ un corps d’exposant caractéristique $p$ tel que $(k:k^p)<+\infty$, $\mathrm B$ une algèbre de séries convergentes sur $k$ et $\mathrm A$ un anneau intègre localisé d’un anneau quotient de $\mathrm B$. Alors $\mathrm A$ est analytiquement réduit.}

Dans le même ordre d’idées, notons la :

\textsc{Proposition 7.} --- \textit{Soit $\mathrm A$ un anneau local nœthérien analytiquement intègre et tel que $\mathrm K_{\widehat{\mathrm A}}$ soit séparable sur $\mathrm K_{\mathrm A}$. Alors la fermeture intégrale de $\mathrm A$ dans toute extension finie $\mathrm L$ de $\mathrm K_{\mathrm A}$ est un $\mathrm A$-module de type fini.}
% slanted
\margin{\handwritten{$\left.\rule{0pt}{2em}\right]$ connu (Nagata)}}

La proposition 7 s’applique aux localités, aux locali\uline{tés analytiques} et aux algèbres de Hopf régulières sur un corps parfait.
\margin{\handwritten{?}}

D’après le lemme 2 (p. 314, chap. 8, \S~13) de (\textsuperscript{11}), on voit alors par exemple qu’une localité analytique vérifiant les conditions du corollaire 2 et normale est analytiquement normale.
% slanted
\margin{\handwritten{$/$ connu}}

(\textsuperscript{*}) Séance du 4 novembre 1963.

(\textsuperscript{1}) \textsc{Berger} et \textsc{Kunz}, \textit{Math., Z.}, \textbf{77}, 1961, p. 314-338.

(\textsuperscript{2}) \textsc{Godement}, \textit{Séminaire Cartan-Chevalley}, 1955-1956, exp. 17, 18, 19.

(\textsuperscript{3}) \textsc{Kunz}, \textit{Math. Z.}, \textbf{76}, 1961, p. 56-74.

(\textsuperscript{4}) \textsc{Lang}, \textit{Ann. Math.}, \textbf{55}, 1952, p. 373-390.

(\textsuperscript{5}) \textsc{Mount}, \textit{Proc. Amer. Math. Soc.} juin 1963.

(\textsuperscript{6}) \textsc{Nagata}, \textit{Locals rings}, Inters., Tracts n\textsuperscript{o} 13, 1962.

(\textsuperscript{7}) \textsc{Nagata}, \textit{On the closedness of singular loci}, Publ. I. H. E. S., n\textsuperscript{o} 2, 1959.

(\textsuperscript{8}) \textsc{Nakai}, \textit{J. Math. Soc. Japan}, \textbf{13}, 1961, p. 63-84.

(\textsuperscript{9}) \textsc{Nakai} et \textsc{Suzuki}, \textit{J. Sc. Hiroshima Univ.}, série A, \textbf{23}, 1960, p. 459-476.

(\textsuperscript{10}) \textsc{Suzuki}, \textit{J. Math. Kyoto Univ.}, \textbf{2}, 1963, p. 157-182.

(\textsuperscript{11}) \textsc{Zariski} et \textsc{Samuel}, \textit{Commutative Algebra}, Van Nostrand.

\hfill(7, \textit{rue Chernoviz, Paris,} 16\textsuperscript{e}).
% archive sheet number: 3
